\documentclass[11pt]{article}
\usepackage[margin=1in]{geometry}
\usepackage{amsmath,amssymb}
\title{ROOT V152: Helicity Representation, Linear Field Equations and Necessary Conditions}
\author{Independent reproducibility audit}
\date{2026-11-25}
\begin{document}
\maketitle
\section*{Units}
Only $c=1$ is adopted, with $\hbar$ and $G$ dimensionful. If $[\hbar]=H$ then
$[G]=L^2/H$, $[R]=L^{-2}$, $[A_\mu]=L^{-1}$, and $[F_{\mu\nu}]=L^{-2}$.
\section*{Geometrical orthogonal frame}
Choose $\mathbf e_1\cdot\mathbf e_2=\mathbf e_1\cdot\mathbf n=\mathbf e_2\cdot\mathbf n=0$ and
$\mathbf e_\pm=(\mathbf e_1\pm i\mathbf e_2)/\sqrt2$. Under an active rotation about $\mathbf n$,
\begin{align}
 R(\varphi)\mathbf e_\pm&=e^{\mp i\varphi}\mathbf e_\pm,\\
 R(\varphi)^{\otimes 2}(\mathbf e_\pm\otimes\mathbf e_\pm)&=e^{\mp 2i\varphi}(\mathbf e_\pm\otimes\mathbf e_\pm).
\end{align}
The same construction produces helicity three and higher: it does not select a physical particle spectrum.
\section*{Spin one}
\begin{align}
 S_A&=-\frac{\hbar}{4g_A^2}\int d^4x\sqrt{-g}\,F_{\mu\nu}F^{\mu\nu},\\
 F_{\mu\nu}&=\partial_\mu A_\nu-\partial_\nu A_\mu,\qquad \nabla_\mu F^{\mu\nu}=0,\\
 f_{\mu\nu}&=i(k_\mu\epsilon_\nu-k_\nu\epsilon_\mu),\quad k^2=0,\quad k\cdot\epsilon=0,\quad \epsilon\sim\epsilon+\lambda k.
\end{align}
\section*{Spin two}
\begin{align}
 S_{\rm GR}&=\int d^4x\sqrt{-g}\frac{R}{16\pi G},\\
 2G^{(1)}_{\mu\nu}&=\partial_\rho\partial_\mu h^\rho{}_{\nu}+\partial_\rho\partial_\nu h^\rho{}_{\mu}-\Box h_{\mu\nu}-\partial_\mu\partial_\nu h
 -\eta_{\mu\nu}(\partial_\rho\partial_\sigma h^{\rho\sigma}-\Box h),\\
 \delta h_{\mu\nu}&=\partial_\mu\xi_\nu+\partial_\nu\xi_\mu,\\
 R^{(1)}_{\mu\nu\rho\sigma}&=-\tfrac12(k_\rho k_\nu h_{\mu\sigma}+k_\sigma k_\mu h_{\nu\rho}-k_\sigma k_\nu h_{\mu\rho}-k_\rho k_\mu h_{\nu\sigma}).
\end{align}
For a transverse-traceless $h_{xx}=H_+$, the curvature-wave amplitude is $R_{0x0x}^{(1)}=\omega^2 H_+/2$.
\section*{Physical limitations}
The orthogonal frame yields rotation representations, but does not generate a free vector action, a metric kinetic term, the gauge group $SU(3)\times SU(2)\times U(1)$, matter chirality, or unique couplings. The equal free-field equations at distinct $g_A$ show that $\alpha$ cannot be inferred from helicity geometry alone. No claim of four-force unification is made.
\end{document}
